% !TeX program = xelatex
\documentclass[11pt,a4paper]{article}
\usepackage[margin=23mm]{geometry}
\usepackage{fontspec}
\setmainfont{Latin Modern Roman}
\usepackage{amsmath,amssymb,mathtools}
\usepackage{xcolor,enumitem,booktabs,longtable,array,xurl,hyperref}
\definecolor{muted}{HTML}{53636C}
\hypersetup{colorlinks=true,urlcolor=blue,linkcolor=blue}
\setlength{\parindent}{0pt}
\setlength{\parskip}{5pt}
\newcommand{\R}{\mathbb R}
\newcommand{\E}{\mathbb E}
\newcommand{\N}{\mathbb N}
\newcommand{\ind}{\mathbf 1}
\newcommand{\Var}{\operatorname{Var}}
\newcommand{\Cov}{\operatorname{Cov}}
\newcommand{\supp}{\operatorname{supp}}
\DeclareMathOperator*{\argmin}{arg\,min}
\DeclareMathOperator*{\argmax}{arg\,max}
\newcommand{\entrymeta}[3]{{\sffamily\small\color{muted}\textbf{\texttt{#1}}\quad|\quad #2\par #3\par}}
\newcommand{\Problemref}[1]{\texttt{#1}}
\newcommand{\JELref}[2]{\texttt{#2} (JEL \texttt{#1})}
\begin{document}
\section*{Supply-chain resilience versus efficiency}
\textbf{Problem ID:} \texttt{L23-120}\quad \textbf{JEL:} \texttt{L23}\par
% BEGIN ECONOMICS STATEMENT
\label{problem:OP-0345}
\label{audited:ECON-080}
{\sffamily\small\color{muted}Original subject group: Trade, globalization and geoeconomics\par}
\label{definitions:problem:30007627}
\textbf{Audit classification:} Published research agenda. The official supply-chain risk priority is verified. The supplier/inventory robust optimization and private-social comparison are editorial, with frontier and dependence assumptions now defined.
\subsection*{Definitions and precise research target}
Consider firms choosing supplier shares \(w\), inventories \(b\) and substitution capabilities \(s\) before uncertain input disruptions. This is a parameterized research agenda: each instance must declare the firms, suppliers, inputs, network and production technology, horizon, contracting and working-capital constraints, and the domains and units of all choices. For each firm-input procurement problem, \(i\) indexes its eligible suppliers; the instance specifies whether shares are quantity or expenditure shares and supplies the total procurement quantity or expenditure and its determination. Inventory units and the production role and domain of substitution capability must also be supplied. The notation \(a=(w,b,s)\) collects these choices over the declared decision-making scope: a single firm with its environment specified, or the economy-wide profile for the social comparison. Feasible configurations belong to the resulting set \(A\); the share simplex alone does not define \(A\).

An instance specifies a measurable space \((\Omega,\mathcal F)\) of primitive shock histories, a nonempty publicly specified set \(\mathcal P\) of probability laws on that space, and the dates and information available for contracting, deliveries, inventory use, substitution and rerouting. The displayed ambiguity set is fixed before the choice and is independent of \(a\). Primitive cost and availability shocks are distinguished from endogenous supplier prices, entry, deliveries, output and consumption, which follow from the declared production and equilibrium rules. A fixed ex ante configuration is distinguished from a contingent rule: if recourse is allowed, its admissible nonanticipative rules and information constraints are part of the instance and its feasible configuration specification. No future shock may be used before it is observed.

Let \(C:A\to\mathbb R\) be normal resource cost and \(L:A\times\Omega\to\mathbb R\) the disruption loss under the declared outcome convention. For every feasible \(a\), require \(L(a,\cdot)\) to be measurable and integrable under every \(P\in\mathcal P\). Write \(R(a)=\sup_{P\in\mathcal P}E_P[L(a,\omega)]\); individual integrability does not imply that \(R(a)\) is finite. The robust social objective and value are
\[J(a)=C(a)+\sup_{P\in\mathcal P}E_P[L(a,\omega)],\qquad J_*=\inf_{a\in A}J(a).\]
The task is to identify the efficiency-resilience frontier and whether private firm choices attain the social optimum when shortages spill through production networks. Estimate switching costs, substitutability and correlated supplier risks rather than assuming suppliers fail independently. Compare diversification, inventory and domestic sourcing using the same loss convention. Include endogenous supplier prices and entry when diversification is adopted widely. Report how optimal decisions vary with upstream position and shock severity, and give bounds where rare-event probabilities are poorly measured. Domestic sourcing is not defined as resilient by assumption; concentrated domestic supply can face common shocks.

\textbf{Additional formal definitions and scope (editorial).} Supplier shares satisfy \(w_i\geq0\), \(\sum_iw_i=1\) for each declared firm-input procurement problem; buffers \(b_i\geq0\) have explicit storage, depreciation and financing accounts; substitution capability \(s\) has an investment cost and an identified production role. Each instance fixes a common normal-output or consumption counterfactual, horizon, numeraire, output valuation, geographic welfare boundary, ownership and social weights. Its accounting decomposition assigns normal real costs to \(C\) and incremental disruption losses to \(L\), including shortage externalities only once. Transfers, financing payments, supplier revenues and rents are distinguished from real resource costs and treated according to the declared welfare boundary and weights. Different baseline or valuation conventions define different instances and must not be mixed in a comparison.

The private sourcing and entry model must specify feasible firm plans, information, beliefs and uncertainty criteria, strategic or price-taking behaviour, entry costs and all market/resource clearing conditions. Its equilibrium correspondence and the outcome rule following each feasible configuration and shock history must be stated. Existence is assumed or established separately. A single-valued \(L(a,\omega)\) requires a declared measurable outcome selection when equilibria are multiple; without such a selection, report separately labelled ranges over permitted equilibria rather than inserting an unspecified equilibrium into \(J\). Evaluate privately induced feasible profiles with the same social \(C\), \(L\) and \(\mathcal P\) used for the planner. A private profile attains the social optimum precisely when its social objective equals \(J_*\); with multiple private equilibria, distinguish whether some or all induced profiles do so. A comparison attributed specifically to shortage externalities holds other primitives, information, beliefs and risk criteria comparable and identifies the uninternalized loss. Comparisons allowing differing beliefs or risk criteria must report those differences rather than attributing the entire gap to externalities. The social disruption term in \(J\) is worst-case expected loss, not expectation under an unspecified single law.

Define the efficiency-resilience frontier as the nondominated attainable finite pairs \((C(a),R(a))\), \(a\in A\): eliminate a pair if another feasible pair is weakly lower in both coordinates and strictly lower in one. This is the attainable frontier, not the frontier of its closure; any unattained limiting pairs are reported separately. Minimizing \(J=C+R\) is one scalar design criterion and does not by itself characterize the entire frontier. Configurations with infinite robust loss are reported separately.

An exact optimizer is conditional on attainment. In particular, for a specified topology on nonempty compact \(A\), lower semicontinuity of the complete objective \(J:A\to\mathbb R\cup\{+\infty\}\) and a finite feasible objective value suffice for a finite attained minimum. Compactness and pointwise finite lower semicontinuous disruption losses alone do not supply these hypotheses; regularity of \(C\), integration and the supremum over laws must also be checked. These are conditional existence assumptions, not restrictions imposed on every instance. For nonempty feasibility and finite \(J_*\), the design target is a feasible choice, or admissible rule where recourse is specified, with \(J(a)\leq J_*+\varepsilon\) for each specified \(\varepsilon>0\). Its abstract existence is distinct from a method for finding it. Report empty feasibility, absence of a finite feasible value, and an infimum of \(-\infty\) separately.

The research deliverables are frontier characterization, private-social comparison and comparative design for declared instances or explicitly delimited model classes. They do not automatically require a finite algorithm or complexity guarantee. A computational specialization must additionally state its input representation, access to laws, costs and equilibrium outcomes, approximation criterion and termination guarantee. Empirical work must distinguish given primitives from estimated quantities and specify observables, sampling and measurement laws, and identification assumptions; where point identification fails, report the compatible identified set. Rare-event bounds must state whether they are model/ambiguity bounds or statistical confidence bounds, their admissible class and assumptions, and any claimed sharpness, informativeness, coverage or uniformity. An empirical protocol or numerical tolerance is not fixed by the broad agenda.
\subsection*{Variants and assumption changes}
\begin{enumerate}
\item \textbf{Known independent disruptions (editorial).} Restrict to a given production technology and independent supplier failures; compare inventory and diversification costs.
\item \textbf{Correlated network disruptions (editorial).} Permit shared geopolitical/natural risks and equilibrium supplier entry/prices; characterize private versus social cost frontiers over an uncertainty set.
\end{enumerate}
\subsection*{References and source audit}
\begin{enumerate}
\item Official January 9, 2026 web version; locator: Interconnected World / Reconfigured trade; 2026 supply-chain priority. Broad agenda; exact model editorial. URL: \url{https://www.bankofengland.co.uk/research/bank-of-england-agenda-for-research}.
\end{enumerate}
\begin{enumerate}\raggedright
\item\hypertarget{definitions:source:30007627:1}{} Bank of England. 2026. \emph{Bank of England Agenda for Research, 2025–2028}. Interconnected World / Reconfigured trade; 2026 supply-chain priority.
\url{https://www.bankofengland.co.uk/research/bank-of-england-agenda-for-research}.
\end{enumerate}

\subsection*{Integrated refinement: ECON-080}
{\sffamily\small\color{muted}Supply-chain resilience versus specialization gains\par}
This refinement adopts the additional source conventions
(page~\pageref{audited:conventions}), subject to its local restrictions.
\textbf{Tail-risk versus ambiguity-based resilience.}
For a declared instance, let \(x=a=(w,b,s)\in A\), including its specified recourse rules if any, and fix a correlated disruption law \(P\) on \((\Omega,\mathcal F)\). Here \(r\) labels the declared regional welfare account: the instance specifies either a particular region or an explicit aggregation over regions with stated weights. The measurable functions \(C_r,L_r:A\times\Omega\to\mathbb R\) are, respectively, its state-dependent resource-cost and incremental disruption-loss accounts, using the same baseline and avoiding double counting. Require \(E_P|C_r(x,\omega)|<\infty\) and \(E_P|L_r(x,\omega)|<\infty\) for every evaluated \(x\). With expectations under this fixed \(P\), consider the additional criterion
\[
J_{\mathrm{tail},r}(x)=\mathbb E[C_r(x,\omega)+L_r(x,\omega)]
       +\lambda\operatorname{CVaR}_\alpha(L_r(x,\cdot)),\qquad
\operatorname{CVaR}_\alpha(L)=\inf_{a\in\mathbb R}
\left\{a+\frac{\mathbb E[(L-a)_+]}{1-\alpha}\right\}.
\]
Here \((z)_+=\max\{z,0\}\), and the scalar \(a\) inside the CVaR definition is only a threshold, unrelated to the configuration. This definition handles atoms at the quantile. Fix \(\alpha\in(0,1)\), \(\lambda\geq0\), common shock scenarios, costs and output valuation. The tail-risk penalty is a preference criterion, not another expected physical loss; \(J_{\mathrm{tail},r}\) is distinct from the original ambiguity-based \(J\). Compare it with worst-case ambiguity and expected-loss criteria on explicitly matched accounts and feasible choices. A reliability comparison must supply a measurable service-failure event \(F_x\), a tolerance \(\delta\in[0,1]\), and whether the constraint is \(P(F_x)\leq\delta\) under a fixed law or \(\sup_{P\in\mathcal P}P(F_x)\leq\delta\). No event or tolerance is implicit. Supplier capacity, inventories, rerouting, substitution and timing constrain feasible diversification. Distinguish private sourcing losses from downstream and distributional social losses. Comparative statements must define the upstream-position statistic and the shock-severity parameter or ordering, state what is held fixed, and report changes in values, frontiers or optimal-choice sets without presuming monotonicity.
\subsection*{Supplied source audit for ECON-080}
Local [S1], [S2], etc. in this audit and the references immediately below
belong to ECON-080, independently of the original entry's reference numbering.
[S1] and [S2] support supply-chain reorientation and resilience. This CVaR/network optimization is an editorial extension, not an exact open question stated by the trade review.
\subsection*{References for integrated refinement ECON-080}
{\footnotesize\raggedright
\par\noindent\textbf{[S1]} David Atkin, Amit Khandelwal, Laura Boudreau, Rafael Dix-Carneiro, Isabela Manelici, Pamela Medina, Brian McCaig, Ameet Morjaria, Luigi Pascali, Heitor Pellegrina, Bob Rijkers, and Meredith Startz (2025). \emph{International Trade}. VoxDevLit 4(2), February 2025.\par \textit{Verified locator / source role:} Primary publisher page checked for review identity, authors, version and background. The specific published agenda is corroborated in [S2]; the exact formal target is editorial.\par \url{https://voxdev.org/voxdevlit/international-trade}\par\smallskip
\par\noindent\textbf{[S2]} Oliver Hanney / VoxDev (living compilation). \emph{Evidence gaps: Key unanswered questions in development economics}. Accessed 8 October 2026. The page currently covers 20 topics; the ``16-topics URL is a historical slug.\par \textit{Verified locator / source role:} Topic heading: International Trade. Supports the broad research agenda; this is not a verbatim quotation or an attribution of the displayed mathematical target.\par \url{https://voxdev.org/topic/evidence-gaps-key-unanswered-questions-across-16-topics-development-economics}\par\smallskip
}

\subsection*{Common conventions: Source mathematical conventions}

These conventions apply to each entry unless explicitly replaced.

\textbf{Models and identification.} A primitive model \(m\) specifies all probability laws, feasible choices, information, equilibrium and measurement rules used by its target. The admissible class \(\mathcal M\) is fixed independently of outcomes. If \(P_O(m)\) is its observable probability law and \(\tau(m)\) the target, its identified set at observable law \(P\) is
\[
 \mathcal I_\tau(P)=\{\tau(m):m\in\mathcal M,\ P_O(m)=P\}.
\]
Point identification means that this set is a singleton. An empty set signals incompatibility of data and assumptions. Coordinate bounds need not describe the joint set. Bounds are sharp only if every retained target is attainable or arbitrarily approachable by compatible models. Finite bounds require explicit restrictions on unobserved quantities; unrestricted confounding, hidden wealth, extrapolation or arbitrary equilibrium selection can yield uninformative sets.

\textbf{Causal interventions.} A potential outcome \(Y_i(p)\) is the outcome under a complete policy regime \(p\), including timing, financing, information and market spillovers. The target population and units are fixed before treatment. With interference, use the complete assignment vector or a justified exposure mapping; individual no-interference notation alone cannot identify an economy-wide intervention. A difference of mean potential outcomes does not require the cross-world joint distribution. Conditioning on post-treatment migration, survival, enrollment or employment changes the estimand and needs separate justification.

\textbf{Statistical statements.} An estimator, confidence set, forecast or test specifies a sampling law, model class and loss/coverage criterion. Uniform \((1-\alpha)\)-coverage means \(\inf_{m\in\mathcal M}\Pr_m\{\tau(m)\in C_n\}\geq1-\alpha\), or an explicitly asymptotic version. The whole parameter space trivially has coverage, so informativeness requires a stated width, risk or power objective. A forecasting improvement is relative to a named benchmark, information set and evaluation population.

\textbf{Equilibrium and optimization.} An equilibrium comprises feasible plans, optimality, beliefs where needed, budget constraints and all market/resource clearing. The primitives or a stated benchmark define each correspondence \(\mathcal E(p;m)\). Nonemptiness is assumed or proved, never inferred just from compact policy sets. A selected-equilibrium target uses a declared selection rule; a robust target quantifies over all permitted equilibria. Use suprema or infima unless attainment is established. An \(\varepsilon\)-optimal policy is feasible and within \(\varepsilon\) of the finite optimum value. Report an empty feasible set separately.

\textbf{Welfare and accounting.} Normative utility, interpersonal normalization, social weights, numeraire, discounting and the use of public receipts are primitives. Behavioral utility need not coincide with the welfare criterion. Household welfare includes ownership income and rents once; adding profits again to compensating variation generally double-counts them. A monetary equivalent is defined by an explicit transfer and price convention, with monotonicity and range conditions. Average effects, mechanism contributions, transfers and welfare losses are different targets. Infinite sums require convergence and dynamic feasible plans require terminal or no-Ponzi conditions.

\textbf{Source versions.} A working-paper date, revision date and journal publication year are distinguished. A DOI for a journal version is not automatically the DOI of an earlier manuscript. Locators refer to the stated version and printed pages where specified. A metadata-only check does not verify a claim about a full-text conclusion. Every entry's source subsection narrows the provenance claim accordingly.

\subsection*{Common conventions: Additional-source status and mathematical conventions}

\label{audited:conventions}

Of the 100 supplied ECON entries, 65 distinct problems appear here; 32 close
matches refine earlier entries, and three duplicate questions map to existing
entries. Appendix~\ref{app:audited-crosswalk} lists every destination.
The attachment's P/R/S/U distinctions and source audits are retained. Its
precise mathematical formalizations are editorial unless the individual audit
says otherwise. Candidate subsidy targets ECON-004--006 retain their U flags;
the resolved approval-core question ECON-007 updates OP-0202 above.
The following supplied conventions govern both this part and the attributed
ECON refinements elsewhere in the catalogue, subject to local restrictions.

\textbf{Parameterized questions.} A problem family lists inputs that must be fixed before an instance is evaluated: population, horizon, policies, information, data law, model class and welfare weights. These are required choices, not quantities the citations are assumed to specify. Undefined applications should not be treated as identified estimands.

\textbf{Indices and integrability.} Sets of agents, firms and regions are finite unless a population measure is explicitly used. \(i,j\) index units, \(r\) regions, \(t\) dates and \(h\) horizons. \(\mathbb E\) and \(\Pr\) refer to the declared population and law; \(\mathbf1\{A\}\) is an indicator. Every displayed expectation and discounted sum needs existence in the stated domain. Infinite horizons use a discount in \((0,1)\) and a convergence condition; finite horizons may use discount one.

\textbf{Optimization and existence.} A displayed \(\max\), \(\min\), or \(\arg\max\) is conditional on attainment. For finite feasible menus this is immediate; general spaces need continuity and compactness/coercivity or another existence argument. Without attainment, the target is the supremum/infimum and arbitrarily near-optimal feasible choices. \(\inf\varnothing=+\infty\). Empty feasibility and an unidentified target are different issues.

\textbf{Potential outcomes.} \(Y_i(z)\) is the outcome under a specified intervention, not the observed conditional mean among people choosing \(z\). In binary comparisons \(\Delta Y_i=Y_i(1)-Y_i(0)\). Specify assignment, treatment versions, compliance, information and observation horizon. Interference is allowed under a full policy vector; compressed exposure notation needs a justified mapping. No interference, consistency and overlap are substantive assumptions, not automatic consequences of notation.

\textbf{Populations and observation.} Fix baseline eligibility and population weights. Conditioning on post-policy employment, survival, relocation or program uptake can change the population being compared. Death is not ordinary loss to follow-up; outcomes truncated by death need a composite convention, joint survival target, or explicitly defined survivor stratum. Fertility-changing interventions need a population functional rather than an unexplained fixed descendant cohort. Measurement and missingness models must be supplied.

\textbf{Identification.} A structural model \(m\in\mathcal M\) implies observable law \(P_m\) and target \(\theta(m)\). Identification means
\[
 P_{m_1}=P_{m_2}\ \Longrightarrow\ \theta(m_1)=\theta(m_2)
 \quad\text{for all }m_1,m_2\in\mathcal M .
\]
Without identification, the target is
\[
 \Theta(P;\mathcal M)=\{\theta(m):m\in\mathcal M,\ P_m=P\}.
\]
Writing \(\theta(P)\) before establishing identification can hide incompatible latent models. Confidence sets additionally need a sampling law, confidence level, asymptotic sequence and uniformity class. Point identification, coverage and informative precision are separate properties.

\textbf{Mediation.} Total effects of randomized assignment are distinct from cross-world natural indirect effects. Belief/effort-channel effects require a meaningful mediator intervention and additional measurement, exclusion or mediation assumptions. Randomization of the initial policy alone does not supply those assumptions. Report bounds or interventional alternatives when the stronger target is unsupported.

\textbf{Equilibrium.} A counterfactual \(W(z)\), \(p(z)\), or \(\mu_t^z\) requires individual optimization, feasibility, government/resource budgets, market clearing and state-transition laws. State equilibrium existence and selection, or report ranges over compatible equilibria. Initial-state integration includes subsequent endogenous transitions; current-state integration must use the policy-dependent distribution. Reduced-form invariance after policy is not assumed.

\textbf{Welfare accounts.} Scores, utility, health, dollars and ecological quantities require explicit conversion before addition. Fees, taxes, subsidies, wages and extortion payments are transfers between parties; distributional weights can make them welfare relevant, but their face value is not automatically a resource loss. State ownership, geographic boundaries, foreign-income weights and fiscal finance. Do not add profits to household welfare already including dividends, or count land capitalization and the underlying benefit twice. A benefit-cost expression must distinguish gross benefits from net welfare and subtract every real cost exactly once.

\textbf{Derivatives and risk.} Log derivatives require positive arguments; local responses require admissible perturbations and specified equilibrium branches. Differentiation under expectations needs regularity/domination, and boundaries may allow only one-sided derivatives. Risk uses a specified law; ambiguity uses a declared nonempty law set on a common history space. Adaptive policies must be nonanticipative. A robust criterion is an editorial choice unless attributed explicitly.

\textbf{Scope overlaps.} Allocation, collective choice, contracts, household bargaining and coordinated policy overlap economics with optimization and game theory. They are retained because they appeared in the original catalogue; no claim of a strictly game-theory-free selection is made.
% END ECONOMICS STATEMENT
\end{document}
